% =====================================================================================
% sections/s6_highdim.tex -- Dimension scaling: PINN versus Deep Ritz for d = 2 to 20.
%
% Every number below is recomputed from the stored per-run records by
%   code/s6_highdim_tables.py  -> data/s6_highdim_numbers.json, data/s6_highdim_tables.tex
% Figures: code/s6_highdim_figs.py (plot-only).
% 16 000 iterations at d = 20: code/s6_highdim_long_d20.py -> data/s6_highdim_long_d20.jsonl.
% Diagnostics of the d = 20 networks: code/s6_highdim_diag.py -> data/s6_highdim_diag_d20.json.
% Sections 6.7 and 6.8: code/s6_highdim_bound.py and code/s6_highdim_remedies.py
%   -> data/s6_highdim_bound_{numbers.json,tables.tex}, data/s6_highdim_remedies_{recheck.json,tables.tex},
%   figures/s6_bound.pdf, figures/s6_remedies_{ratios,altridge}.pdf, from the result files of
%   research_highdim_bound/ and research_highdim_remedies/ (no training).
% Table rows between GENERATED markers are written by these scripts (code/inject.py finds the
% marker in any section file).  Further results, tables and figures: sections/S5_highdim.tex.
% =====================================================================================
\section{Dimension scaling: PINN versus Deep Ritz for \texorpdfstring{$d=2$}{d = 2} to 20}%
\label{s6_highdim:sec}

Neural solvers are usually motivated by problems in many dimensions, where a mesh is out
of reach~\citep{eyu2018,sirignano2018,han2018}. This section measures how the accuracy and
the cost of the PINN and of the Deep Ritz method change with the dimension, on the Laplace
problem \Pone~\eqref{s2_problems:eq:p1} and the Poisson problem \Ptwo~\eqref{s2_problems:eq:p2}
for $d=2,3,5,10,20$, with the protocol of Section~\ref{s3_methods:sec:highdim}: one network
($d$--64--64--64--1, tanh), 1024 interior and 1024 boundary points redrawn at every
iteration, \Adam\ with two learning-rate reductions, 4000 iterations unless stated,
$\lam=1000$ for the PINN, $\bet=100$ (\Pone) and $\bet=1$ (\Ptwo) for Deep Ritz, and the
error of the final network on 50\,000 held-out points. Each cell has three seeds (0, 1, 2).
% src: research_highdim/scripts/hd_core.py (train: torch.manual_seed(seed), generator 1000+seed for both methods)
Within a seed the two methods start from the same weights and see the same points, so the
comparison is paired. One fact frames the results: \Ptwo\ has $\dn\uex=0$ and is the best
case for the penalised energy, which is then unbiased, whereas \Pone\ is the generic,
biased case (Proposition~\ref{s3_methods:prop:robin}); the PINN loss has no such bias
(Remark~\ref{s3_methods:rem:consistent}). Sections~\ref{s6_highdim:sec:equal-its}
to~\ref{s6_highdim:sec:weights} report the comparison. Section~\ref{s6_highdim:sec:bound} then
proves an explicit stability constant for the unit cube, which turns the two terms of the PINN
loss into a computable bound on the error of any network, and Section~\ref{s6_highdim:sec:remedies}
tests two cheap remedies for the plateau that both methods reach on \Ptwo\ at $d=20$.

% -------------------------------------------------------------------------------------
\subsection{Accuracy at equal iterations}\label{s6_highdim:sec:equal-its}

\begin{figure}[tbp]
\centering
\includegraphics[width=0.8\textwidth]{s6_highdim_error.pdf}
\caption{Held-out errors against the dimension for \Pone\ (left) and \Ptwo\ (right):
relative $L^2$ error $\errL$ (top) and relative $H^1$-seminorm error $\errH$ (bottom) on
50\,000 test points. Filled symbols joined by lines: mean over seeds 0--2 after 4000
iterations (the values of Tables~\ref{s6_highdim:tab:p1} and~\ref{s6_highdim:tab:p2});
small dots: the individual seeds. Open symbols: the same after 16\,000 iterations, run for
\Pone\ at $d=10$ and for both problems at $d=20$. Dashed line: $\errL$ of the best affine
approximation of $\uex$ (Proposition~\ref{s2_problems:prop:floors}). PINN in orange
(squares), Deep Ritz in blue (circles); the two methods are offset horizontally for
legibility.}
\label{s6_highdim:fig:error}
% src: research_highdim/results/runs*.jsonl; data/s6_highdim_long_d20.jsonl; data/closed_forms.json; script code/s6_highdim_figs.py
\end{figure}

% src: research_highdim/results/summary.csv; research_highdim/results/derived.txt (ratios, range overlap; boundary vs interior)
For every $d\le10$ the PINN is the more accurate method on both problems
(Tables~\ref{s6_highdim:tab:p1} and~\ref{s6_highdim:tab:p2},
Figure~\ref{s6_highdim:fig:error}): the ratio of the mean errors lies between 1.46 and 4.29,
and in each of the eight cells the three PINN seeds are below the three Deep Ritz seeds. The
$H^1$-seminorm errors are ordered in the same way. At $d=20$ the two methods are level on
\Pone\ (ratio 0.99, overlapping ranges) and Deep Ritz is 12 per cent lower on \Ptwo. In every
cell the relative $L^2$ error on 20\,000 held-out boundary points is between 0.65 and 2.11
times the interior error.
% src: verify_research_highdim/results/v_summary.txt (rows 'main'); data/s6_highdim_numbers.json (vdim_main, vdim_P1_d20_ritz_sd_5_seeds = 4.06e-3)
The verification code \textsf{V-dim} confirms the ordering: at $d=2$ and $d=10$, on both
problems, each of its PINN seeds is below each of its Deep Ritz seeds
(Table~\ref{s6_highdim:tab:vdim}). It also shows that three seeds understate the spread (on
\Pone\ at $d=20$ its five Deep Ritz seeds have a standard deviation of $\sci{4.1}{-3}$ against
$\sci{1.8}{-3}$ for our three), so we rely only on differences that hold for every seed.

% -------------------------------------------------------------------------------------
\subsection{Growth with the dimension, and trivial-predictor floors}\label{s6_highdim:sec:growth}

% src: research_highdim/results/derived.txt (log-log slopes of the geometric-mean error); recomputed in data/s6_highdim_numbers.json (slopes_geomean); research_highdim/results/summary.csv (laplace, ritz: rel_l2_mean 1.8821e-2, 1.4214e-2, 1.1126e-2, 1.8964e-2 at d = 2, 3, 5, 10)
At this fixed budget the error grows with $d$ for both methods, except Deep Ritz on \Pone\ for
$d\le10$, where it is not monotone ($\sci{1.88}{-2}$, $\sci{1.42}{-2}$, $\sci{1.11}{-2}$ and
$\sci{1.90}{-2}$ at $d=2$, 3, 5 and 10). A least-squares line through
the logarithm of the geometric-mean error against $\log d$ over $d=2,\dots,10$ has slope 0.48
(PINN) and 0.01 (Deep Ritz) on \Pone, and 1.71 and 1.46 on \Ptwo; with $d=20$ included the
slopes are 0.86, 0.41, 1.92 and 1.35. These numbers describe one budget and four or five
dimensions. They are not rates: Section~\ref{s6_highdim:sec:budget} shows that they change
with the budget.
% src: data/closed_forms.json (floors); research_highdim/results/summary.csv (rel_l2_centered_mean)
Relative errors must be read against the floors of Proposition~\ref{s2_problems:prop:floors}.
The solution of \Pone\ concentrates around its mean as $d$ grows, so both floors fall with $d$
and plain $\errL$ flatters \Pone\ in high dimension: the centred error of the PINN rises from
$\sci{9.2}{-3}$ at $d=2$ to $\sci{3.4}{-2}$ at $d=10$ and $\sci{1.73}{-1}$ at $d=20$, so the
``4.7 per cent'' of Table~\ref{s6_highdim:tab:p1} at $d=20$ is 17 per cent of the fluctuation
of $\uex$. For \Ptwo\ the affine floor is $0.1203$ in every dimension, and at $d=20$ both
methods end above it ($0.144$ and $0.126$): after 4000 iterations neither improves on the
best affine function.

% -------------------------------------------------------------------------------------
\subsection{The networks at \texorpdfstring{$d=20$}{d = 20}, and the effect of the budget}%
\label{s6_highdim:sec:budget}

% src: verify_research_highdim/results/v_affine.txt; verify_research_highdim/results/v_resid.txt (the 4000-iteration networks, three seeds each); data/s6_highdim_diag_d20.json (re-computation on other points: 0.18-0.20, 0.90; 0.82-0.83, 0.13-0.16, 0.35-0.36)
The two methods are above the affine floor on \Ptwo\ at $d=20$ for different reasons. Split a
function into its least-squares affine fit and a remainder. For the Deep Ritz networks the
remainder has 0.18 to 0.20 times the norm of the remainder of $\uex$, and the residual
$\norm{\lap\unet+f}/\norm{f}$ is 0.90: these networks are still essentially affine. The PINN
networks are not. Their residual is 0.35 to 0.37, and their remainder has 0.83 to 0.84 times
the right norm, but its correlation with the remainder of $\uex$ is only 0.14 to 0.16: the
PINN has built curvature of the right size and the wrong shape. The following elementary
bound says where its error sits.

\begin{remark}\label{s6_highdim:rem:split}
Let $v\in C^2(\overline{\Om})$, $\Om=(0,1)^d$, $e=v-\uex$ and $r=\lap v+f=\lap e$. Let
$e_0\in H^1_0(\Om)$ be the weak solution of $\lap e_0=r$; then $e-e_0$ is harmonic and has
the boundary values of $e$. Expanding $e_0$ and $r$ in the orthonormal Dirichlet
eigenfunctions $\prod_{i}\sqrt2\sin(k_i\pi x_i)$ of $-\lap$, whose eigenvalues
$\pi^2\sum_i k_i^2$ are at least $d\pi^2$, gives
$\norm{e_0}_{L^2(\Om)}\le\norm{r}_{L^2(\Om)}/(d\pi^2)$. For \Ptwo, where
$\norm{f}=\pi^2\norm{\uex}$, this reads
$\norm{e_0}/\norm{\uex}\le d^{-1}\norm{\lap v+f}/\norm{f}$.
\end{remark}

\noindent
% src: data/s6_highdim_numbers.json (e0_bound_P2_d20); research_highdim/results/summary.csv (bd_rel_l2_mean at d = 20: 0.158, 0.133)
With the measured residuals, the part of the error that the interior residual accounts for is
at most $0.37/20<0.02$ of $\norm{\uex}$ for the PINN and $0.90/20=0.045$ for Deep Ritz, out of
total errors of $0.144$ and $0.126$. For both methods the error is thus mostly a harmonic
function fixed by the mismatch on the boundary, which is 16 per cent (PINN) and 13 per cent
(Deep Ritz) in relative $L^2$ norm. Corollary~\ref{s6_highdim:cor:bound}(c) completes this
remark with an explicit constant for the harmonic part.

\begin{table}[tbp]
\centering
\caption{Effect of the budget: $\errL$ in units of $10^{-2}$, mean $\pm$ sample standard
deviation over seeds 0--2. ``16\,000'' is the same protocol with four times as many
iterations (learning-rate reductions at 8000 and 12\,000). ``Equal compute'' is Deep Ritz
with the number of iterations given in brackets, chosen to match the compute of the
4000-iteration PINN run (Section~\ref{s6_highdim:sec:compute}). Last column: affine floor.
A dash means not run.}
\label{s6_highdim:tab:budget}
\small
\setlength{\tabcolsep}{4pt}
% src: research_highdim/results/summary.csv (4000); research_highdim/results/extensions.json (d = 10: 16 000 and equal compute)
% src: data/s6_highdim_long_d20.jsonl (d = 20: 16 000; script code/s6_highdim_long_d20.py)
% generated by code/s6_highdim_tables.py -> data/s6_highdim_tables.tex
\begin{tabular}{@{}lr cc cc c c@{}}
\toprule
 & & \multicolumn{2}{c}{4000 iterations} & \multicolumn{2}{c}{16\,000 iterations} & equal compute & \\
\cmidrule(lr){3-4}\cmidrule(lr){5-6}\cmidrule(lr){7-7}
 & $d$ & PINN & Deep Ritz & PINN & Deep Ritz & Deep Ritz & floor\\
\midrule
\Pone & 10 & $1.260\pm0.084$ & $1.896\pm0.159$ & $0.588\pm0.027$ & $0.728\pm0.039$ & $0.934\pm0.047$ (10\,000) & 13.9 \\
\Ptwo & 10 & $2.168\pm0.482$ & $7.488\pm2.553$ & -- & -- & $1.376\pm0.111$ (9000) & 12.0 \\
\Pone & 20 & $4.663\pm0.248$ & $4.613\pm0.178$ & $0.757\pm0.023$ & $0.870\pm0.038$ & -- & 10.2 \\
\Ptwo & 20 & $14.424\pm0.176$ & $12.633\pm0.022$ & $5.832\pm0.858$ & $9.102\pm0.808$ & -- & 12.0 \\
\bottomrule
\end{tabular}
\end{table}

% src: research_highdim/results/derived.txt (plateau block); data/s6_highdim_numbers.json (plateau_val_mean)
The validation curves (Figure~\ref{s6_highdim:fig:curves}) show that these errors are limited
by the budget. Within 250 iterations the error falls to about the affine floor; it then stalls
before descending, and the stall lengthens with $d$; at $d=20$ neither method has left the
plateau on \Ptwo\ when the 4000-iteration run ends. That networks fit the smoothest component
first is a known tendency~\citep{rahaman2019}; we cite it as context, not as an explanation.
% src: research_highdim/results/extensions.json (d = 10); data/s6_highdim_long_d20.jsonl, data/s6_highdim_numbers.json (extensions, budget_ratios, growth_P1_d10_to_d20)
With four times the budget (Table~\ref{s6_highdim:tab:budget}) the errors at $d=10$ on \Pone\
fall by a factor 2.1 (PINN) and 2.6 (Deep Ritz), and at $d=20$ on \Pone\ by a factor 6.2 and
5.3. On \Ptwo\ at $d=20$ both methods leave the plateau and all six runs end below the affine
floor: $\sci{5.83}{-2}$ for the PINN and $\sci{9.10}{-2}$ for Deep Ritz. These runs are not
converged either: the error is still falling when the learning rate is reduced. With the larger
budget the PINN is again the more accurate method in every seed at $d=20$ (ratio 1.15 on
\Pone, 1.56 on \Ptwo), so the level or reversed ranking after 4000 iterations was an effect of
the budget and, at least in part, of the uncentred input:
% src: research_highdim_remedies/results/summary_exploratory.json (P1/d20/pinn/lift1/4000: 0.0219 to 0.0317; P1/d20/ritz/lift1/4000: 0.0442 to 0.0551)
with the input centred, the PINN is the more accurate method on \Pone\ at $d=20$ after 4000
iterations as well, in every seed (Section~\ref{s6_highdim:sec:remedies}). Much of the growth with
$d$ was also an effect of the budget: from $d=10$ to $d=20$ the PINN error on \Pone\
grows by a factor 3.7 after 4000 iterations and 1.3 after 16\,000 (Deep Ritz: 2.4 and 1.2).
For \Pone, Table~\ref{s6_highdim:tab:p1} is therefore to be read as ``the number of iterations
needed for a given accuracy grows with $d$''. For \Ptwo\ the corresponding long run at $d=10$
was not made, and for neither problem was the approximation power of the network tested
separately.

% -------------------------------------------------------------------------------------
\subsection{Cost per iteration}\label{s6_highdim:sec:cost}

Deep Ritz needs only $\grad\unet$, one reverse pass whatever the dimension. The PINN needs
$\lap\unet$, that is $d$ second derivatives, either by one further reverse pass per coordinate
or by the forward propagation of Section~\ref{s3_methods:sec:highdim}, which evaluates the same
loss. Two separate measurements on the same shared machine (Table~\ref{s6_highdim:tab:cost},
Figure~\ref{s6_highdim:fig:cost}) agree in the following.
% src: data/s6_highdim_numbers.json (cost_growth_d2_to_d100, nested_over_fwd_d_ge_10, cost_primary, cost_verifier); research_highdim/results/cost_vs_d.csv; verify_research_highdim/results/v_bench.txt
From $d=2$ to $d=100$ the cost of a Deep Ritz iteration grows by a factor 1.98 or 1.55, from
3.3 to 5.1\,ms in the quieter measurement, and that of a PINN iteration roughly linearly, as is
known~\citep{hu2024}: by a factor 13.9 or 12.7 (forward Laplacian) and 32.5 or 26.2 (nested).
At $d=100$ a PINN iteration costs about 11 to 13 times a Deep Ritz iteration with the forward
Laplacian and about 29 to 32 times with the nested one; for $d\ge10$ the forward Laplacian is
1.8 to 2.9 times cheaper than the nested one. The milliseconds themselves depend on the load
and are not portable.

% -------------------------------------------------------------------------------------
\subsection{Accuracy at equal compute}\label{s6_highdim:sec:compute}

Because a Deep Ritz iteration is cheaper, a comparison at equal iterations understates the
method. Deep Ritz was therefore also run with the compute of the 4000-iteration PINN run
at $d=10$.
% src: research_highdim/scripts/run_study.py (phase ext, rule); data/s6_highdim_numbers.json (equal_compute_rule: 9594 -> 10000, 8889 -> 9000)
The number of iterations was fixed before the runs, as the mean CPU time of the three PINN runs
divided by the mean CPU time per iteration of the three Deep Ritz runs, rounded to a thousand:
10\,000 on \Pone\ and 9000 on \Ptwo, that is 2.5 and 2.25 Deep Ritz iterations for one PINN
iteration (the cost ratio measured at $d=10$ is 2.23 and 2.48 in the two measurements).
% src: research_highdim/results/extensions.json (phase equaltime); research_highdim/results/summary.csv; data/s6_highdim_numbers.json (equal_compute_gain); verify_research_highdim/results/v_summary.txt (rows 'equaltime', 'eq9k'); notes/VERIFICATION.md (sec. 3.1, V10; sec. 3.3, item 5)
The ranking reverses (Table~\ref{s6_highdim:tab:budget}): Deep Ritz reaches
$\sci{9.34}{-3}\pm\sci{4.7}{-4}$ against $\sci{1.26}{-2}\pm\sci{8.4}{-4}$ for the PINN on
\Pone, and $\sci{1.38}{-2}\pm\sci{1.1}{-3}$ against $\sci{2.17}{-2}\pm\sci{4.8}{-3}$ on \Ptwo,
with disjoint seed ranges in both cases. \textsf{V-dim} agrees (seeds 7--9:
$\sci{8.74}{-3}$, $\sci{9.09}{-3}$, $\sci{9.02}{-3}$ on \Pone; $\sci{1.38}{-2}$,
$\sci{1.24}{-2}$, $\sci{1.50}{-2}$ on \Ptwo). In the re-check of
Section~\ref{s6_highdim:sec:bound}, with a separately written trainer and three seeds, Deep Ritz
at the CPU time of the PINN was also the more accurate on \Pone\ at $d=20$ (descriptive).
% src: data/s6_highdim_bound_numbers.json (recheck_r2: P1_ritz_eqcpu_d20 rel_err 0.0092-0.0104, P1_pinn_d20 0.0459-0.0484)
The ranking of the two methods thus depends on whether iterations or compute are counted. It
does not show that Deep Ritz is the better method: at 16\,000 iterations each the PINN is ahead
at $d=10$ on \Pone\ and at $d=20$ on both problems, although on \Pone\ that lead is of the size
of the penalty bias of Deep Ritz (Section~\ref{s6_highdim:sec:weights}).

% -------------------------------------------------------------------------------------
\subsection{Sensitivity to the penalty weights, and the penalty bias}\label{s6_highdim:sec:weights}

% src: research_highdim/results/sweep.csv (val_rel_l2; d = 5, seed 100); data/s6_highdim_robin_bias.json (weight_sweep_d5: bias 0.1640, 0.02645, 2.991e-3, 3.047e-4; trained_over_bias 0.98, 0.88, 3.44, 36.0)
The weights were selected at $d=5$ from one run per value (Table~\ref{s6_highdim:tab:sweep}).
On \Pone\ the error of Deep Ritz falls by a factor 16 from $\bet=1$ to $\bet=100$ and is level
beyond; the computed penalty bias at $d=5$ explains the first part (it is 0.164 for $\bet=1$
and $\sci{2.64}{-2}$ for $\bet=10$, and these two runs end at 0.98 and 0.88 times it), while
for $\bet=100$ and $1000$ the runs end 3.4 and 36 times above the bias, at an error limited by
the optimisation. On \Ptwo, Deep Ritz is flat, as expected from $\dn\uex=0$, and the PINN needs
$\lam\ge10$.
% src: verify_research_highdim/results/v_summary.txt (rows 'beta', 'lam1e4', 'main' at d = 10); notes/VERIFICATION.md (sec. 3.3, item 3: V11, V13)
These are statements about $d=5$. At $d=10$ larger weights than the chosen ones do not help:
on \Ptwo\ (\textsf{V-dim}, seeds 7--9) Deep Ritz with $\bet=10$ or $100$ and the PINN with
$\lam=10^4$ stay on the affine plateau for the whole run (Supplementary
Section~\ref{S5:sec:weights}).

% src: data/s6_highdim_robin_bias.json (series: beta = 100; against_trained_networks); research_highdim_bound/results/bias_table.json (exact_over_article_series 0.99957 to 1.00026)
The penalty bias of Deep Ritz on \Pone\ was computed for every $d$ of the study by a series in
the Robin eigenfunctions and, exactly, up to $d=100$ by the representation of
Proposition~\ref{sA_proofs:prop:exact};
the two agree to $\sci{4.3}{-4}$ relative (Supplementary Section~\ref{S5:sec:bias} and
Table~\ref{s6_highdim:tab:biasexact}). At $\bet=100$ the relative bias $\norm{\ubeta-\uex}/\norm{\uex}$
falls from $\sci{5.5}{-3}$ at $d=2$ to $\sci{7.4}{-4}$ at $d=20$, roughly like $0.013/d$ to
$0.015/d$ for $d\ge3$. At $d=2$ it is about the PINN's whole error ($\sci{6.1}{-3}$) and 0.29 of
the Deep Ritz error. After 4000 iterations it is 27 to 33 per cent of the Deep Ritz error for
$d\le5$ and 8 and 2 per cent at $d=10$ and 20, so at that budget the gap between the methods is
mostly optimisation. After 16\,000 iterations it is 1.05 and 0.65 times the gap to the PINN at
$d=10$ and 20: the lead of the PINN on \Pone\ at the larger budget is of the size of the bias
that the fixed $\bet=100$ imposes on Deep Ritz, and it does not show that the PINN optimises
its loss better.

% src: research_highdim_bound/results/bias_table.tex (bound(b)/exact and article/exact at beta = 100); research_highdim_bound/results/bias_table.json
The representation of Proposition~\ref{sA_proofs:prop:exact}, a one-dimensional integral over
products of one-dimensional Robin eigenfunction sums, is valid for every $d$ and $\bet$. Against
these exact values, the bound of Corollary~\ref{s6_highdim:cor:bound}(b) of the next section is 1.71 to 4.43 times the exact bias at
$\bet=100$ for $d=2,\dots,100$, against 6.5 to 49.6 times for the bound of
Proposition~\ref{s3_methods:prop:robin}(c). The numbers suggest that $d\,c(d)\to\sqrt{8/3}$ for
$c(d)=\lim_{\bet\to\infty}\bet\norm{\ubeta-\uex}/\norm{\uex}$
(Conjecture~\ref{s6_highdim:conj:dc}).


% -------------------------------------------------------------------------------------
\subsection{An explicit stability constant and a computable error bound}%
\label{s6_highdim:sec:bound}

Remark~\ref{s6_highdim:rem:split} bounds the part of the error that the interior residual
accounts for; the rest is harmonic and has the boundary values of $v-g$. Its size is governed by
\begin{equation}\label{s6_highdim:eq:kappa}
\kappa_d=\sup\Bigl\{\frac{\nbd{\dn\psi}}{\norm{\lap\psi}}:\ \psi\in H^2(\Om)\cap H^1_0(\Om),\
\psi\ne0\Bigr\},\qquad\Om=(0,1)^d,
\end{equation}
where $\nbd{\cdot}$ is the norm of $L^2(\dOm)$ with surface measure ($\abs{\dOm}=2d$). By
Green's formula, $-\int_\Om w\,\lap\psi=-\int_{\dOm}w\,\dn\psi$ for harmonic $w$ and $\psi=0$ on
$\dOm$, so $\kappa_d$ is the norm of the harmonic extension from $L^2(\dOm)$ to $L^2(\Om)$
(Supplementary Section~\ref{sA_proofs:sec:kappa}). This is a classical quantity:
$\kappa_d^2=1/\delta_1$, where $\delta_1$ is the first eigenvalue of the biharmonic Steklov
problem $\lap^2u=0$ in $\Om$, $u=0$ and $\lap u=\delta_1\dn u$ on
$\dOm$~\citep{ferrero2005fourth,bucur2009first,bucur2011first}
(Proposition~\ref{s6_highdim:prop:steklov}). For convex domains of minimal width $\ell$, Payne's
bound $\delta_1\ge2/\ell$~\citep{payne1968bounds}, as stated in~\citet[Lemma~5.3]{bucur2011first}
(see also~\citep{philippin2010extending}), gives
$\kappa_d^2\le\frac12$ on the unit cube in every dimension. The following theorem gives the exact
order in $d$.
% src: refs.bib (bucur2011first, annote: Lemma 5.3, Payne's bound; Section 3, Fichera's duality)

\begin{theorem}[The harmonic-extension constant of the unit cube]\label{s6_highdim:thm:kappa}
\StmtKappa
\end{theorem}

\begin{proposition}[Biharmonic Steklov eigenvalue]\label{s6_highdim:prop:steklov}
\StmtSteklov
\end{proposition}

\begin{corollary}[Harmonic part, penalty bias and a computable error bound]\label{s6_highdim:cor:bound}
\StmtCertificate
\end{corollary}

\noindent The proofs are in Supplementary Section~\ref{sA_proofs:sec:kappa}; every step was also
checked numerically (finite differences on the faces, the partial-fraction identity, Rayleigh
quotients of the extremal sequences, and the bound on 200 harmonic test functions per dimension).
% src: research_highdim_bound/results/theorem2_checks.txt (A1-A7); research_highdim_bound/results/corollary_checks.txt; research_highdim_bound/results/posthoc_analyses.json (step8_check)
Corollary~\ref{s6_highdim:cor:bound}(c) is Remark~\ref{s6_highdim:rem:split} completed by an
explicit constant for the harmonic part; it needs neither $\uex$ nor a mesh, only the two
quantities that the PINN loss estimates. The norm $\nbd{\cdot}$ uses the surface measure, so in
terms of the root-mean-square boundary misfit, the quantity that the boundary term of the loss
estimates, the factor is $\sqrt{2dU_d}$, which grows like $d^{1/4}$ (1.08 at $d=2$, 1.71 at
$d=20$, 2.53 at $d=100$; Table~\ref{s6_highdim:tab:kappa}): in high dimension the boundary
misfit weighs more in the bound, not less.
% src: research_highdim_bound/results/table_theorem2.tex (columns sqrt(U_d), sqrt(2dU_d) and 2 sqrt(1+1/(d pi^2)))
Part~(b) replaces the constant $2C_d$ of Proposition~\ref{s3_methods:prop:robin}(c), which lies
between 2.00 and 2.10 for all $d$, by $\sqrt{U_d}$, which is 0.540, 0.326, 0.270 and 0.179 at
$d=2$, 10, 20 and 100.

\paragraph{Comparison with the classical bounds.} Most of that improvement is classical: Payne's
bound already gives the constant $\sqrt{1/2}$, a factor 2.83 to 2.97 below $2C_d$ for
$d=1,\dots,100$. Against Payne's constant, Theorem~\ref{s6_highdim:thm:kappa} gains the factor
$\sqrt{(1/2)/U_d}$, which is 1 at $d=1$, 1.31 at $d=2$, 2.17 at $d=10$ and 3.95 at $d=100$.
% src: research_highdim_bound/results/posthoc_analyses.json (payne_compare: article_over_payne 2.968 at d = 1 to 2.830 at d = 100; sqrt_payne_over_sqrtU); research_highdim_bound/results/table_payne.tex
The elementary bound through the torsion function, $\kappa_d^2\le m_d$ with $m_d$ the largest
normal derivative of the torsion function of the cube on its boundary
(Remark~\ref{sA_proofs:rem:steklov}), already decreases with $d$, but slowly: $m_d=0.338$, 0.195, 0.170 and 0.137 at
$d=2$, 10, 20 and 100. Against it the theorem gains $\sqrt{m_d/U_d}=1.08$, 1.36, 1.53 and 2.07
(Table~\ref{s6_highdim:tab:kappa}). What the theorem adds is therefore the exact order
$d^{-1/2}$ of $\kappa_d^2$, not a large constant at small $d$.
% src: data/s6_highdim_torsion.json (m_d and sqrt(m_d/U_d); code/s6_highdim_torsion.py, two quadratures agreeing to 8e-11)
The elementary lower bound $1/(2d)$, from the constant harmonic function, is larger than $L_d$ for
$2\le d\le8$; $L_d$ is what gives the order $d^{-1/2}$ from below. How sharp $U_d$ is remains
open: the computed Rayleigh quotients fall from $0.985\,U_d$ at $d=2$ to $0.874\,U_d$ at $d=10$
and $0.637\,U_d$ at $d=100$, but the subspaces that can be afforded shrink with $d$, so these
numbers neither support nor refute $\kappa_d^2/U_d\to1$; what is proved is
$\frac12(1+o(1))\le\kappa_d^2/U_d\le1$. For the unit square see Supplementary
Section~\ref{S5:sec:square}.
% src: research_highdim_bound/check/out_r2/r2_thm2_edges.json (L_8 = 0.06037 < 1/16, L_9 = 0.05639 > 1/18); research_highdim_bound/results/table_theorem2.tex; research_highdim_bound/results/best_lower_bounds.json

\paragraph{The bound on the trained networks.} Corollary~\ref{s6_highdim:cor:bound}(c) was
evaluated on all 88 saved networks of this section, with norms estimated by Monte Carlo on
$2\times10^5$ new interior and $2\times10^5$ new boundary points per $d$ and, as a second set, on
the held-out test points. Three hypotheses were pre-registered and frozen before the
confirmatory evaluation: H1, the bound exceeds the error of every network; H2, its efficiency
$\eta=(B_{\mathrm{int}}+B_{\mathrm{bd}})/\norm{v-\uex}$ is at most 3 for every confirmatory
network with $d\ge5$; H3, the seed-mean share $B_{\mathrm{bd}}/(B_{\mathrm{int}}+B_{\mathrm{bd}})$
does not decrease with $d$. Seven of the 88 networks had been seen in a pilot evaluation and are
treated as non-confirmatory (Supplementary Section~\ref{S5:sec:bound}).
% src: research_highdim_bound/PREREG_T2.md; research_highdim_bound/results/inventory.json (88 checkpoints, 7 identical to pilot networks)
% src: research_highdim_bound/results/decisions.json (H1: violations [], min_z_D fresh 222.06, hd 75.49, min_eta 1.744; H2: n_tested 57, 21 failures, eta_range [1.744, 5.456]; H3); research_highdim_bound/results/posthoc_analyses.json (eta_by_d_all88: d = 2 max 24.82, d = 3 max 8.31)
H1 holds: on both sets and for all 88 networks the bound is above the error, by at least 75
standard errors, and the smallest efficiency is 1.744. H2 fails: $\eta$ exceeds 3 for 21 of the
57 confirmatory networks with $d\ge5$, with a maximum of 5.46 (Deep Ritz on \Ptwo\ at $d=5$);
below $d=5$, where H2 made no claim, $\eta$ is up to 8.31 at $d=3$ and up to 24.8 at $d=2$. The
bound is therefore a valid but loose certificate at small $d$, not an error estimator. H3 holds
for all four (problem, method) pairs: the seed-mean share rises from 0.19 to 0.29 at $d=2$
(0.056 for Deep Ritz on \Ptwo) to 0.84 to 0.98 at $d=20$ (Figure~\ref{s6_highdim:fig:bound}). Part
of this trend is built into the constants, $1/(d\pi^2)$ for the interior and
$\sqrt{2dU_d}\sim d^{1/4}$ for the boundary, so H3 alone does not show that the error itself is
dominated by the boundary misfit; at $d=20$ Remark~\ref{s6_highdim:rem:split} does show it for the
networks tested, since $B_{\mathrm{int}}$ is at most 0.36 of the error on every main network.
% src: research_highdim_bound/results/share_by_d.csv; data/s6_highdim_bound_numbers.json (B_int_over_error_main_networks: d = 20 at most 0.355, d = 10 up to 0.805)
% src: research_highdim_bound/check/out/c5_compare.json (eta_rel_diff_range [-0.0078, 0.0061]); research_highdim_bound/check/out/c4_fresh.jsonl (24 rows); research_highdim_bound/check/out_r2/r2_fresh.jsonl (21 rows); data/s6_highdim_bound_numbers.json (n_networks_bound_valid: 88 + 24 + 21; eta_d10_d20_all_networks 1.916 to 3.444; eta_by_d_all_networks d = 2: max 24.82; recheck_r2: pairs_ranked_like_error_equal_cpu 9/9, equal_iterations 4/9)
The re-check repeated the evaluation with its own evaluator ($\eta$ within $-0.78$ to $+0.61$ per
cent on all 88 networks, the same decisions) and retrained 45 networks with separately written
trainers and other seeds, on \Pone, \Ptwo\ and two problems outside the study; this part was not
pre-registered. Over all 133 networks evaluated the bound held in every case; its efficiency is
1.92 to 3.44 at $d=10$ and 20 and up to 24.8 at $d=2$. At $d=20$ and equal CPU time the bound
ranks all $3\times3$ pairs of a retrained PINN and Deep Ritz network as the error does; at equal
iterations, where the errors of the two methods lie within the seed spread, it ranks only 4 of
the 9 pairs correctly: a bound within a factor of about 2 of the error is not meant to resolve
error differences much smaller than that factor.

\begin{figure}[tbp]
\centering
\includegraphics[width=0.78\textwidth]{s6_bound.pdf}
\caption{(a) The bounds $U_d$, $L_d$ and $1/(2d)$ of Theorem~\ref{s6_highdim:thm:kappa} on
$\kappa_d^2$, Payne's bound $\frac12$, the torsion-function bound $m_d$ of the cube
(Remark~\ref{sA_proofs:rem:steklov}), the largest computed Rayleigh quotients, and the square of the
constant $2C_d$ of Proposition~\ref{s3_methods:prop:robin}(c). (b) Efficiency $\eta$ and (c) boundary
share of the bound of Corollary~\ref{s6_highdim:cor:bound}(c) for all 88 saved networks, new
Monte-Carlo set: the main runs (lines: seed means; open markers: the seven networks seen in the
pilot), the 16\,000-iteration and equal-compute runs at $d=10$ (triangles), and the 16 networks of
the weight sweep at $d=5$ (grey, drawn slightly to the left of $d=5$). The dotted line in (b) is the
pre-registered limit 3.}
\label{s6_highdim:fig:bound}
% src: research_highdim_bound/results/theorem2_checks.json; research_highdim_bound/results/restricted_lanczos.json; research_highdim_bound/results/eval_confirm.jsonl (88 networks: main 60, long and equaltime 12, sweep 16); data/s6_highdim_torsion.json; script code/s6_highdim_bound.py (plot only)
\end{figure}

\paragraph{Relation to the literature.} Besides the classical facts recalled above (Fichera's
duality, Payne's bound, the torsion function), numerical bounds for $\delta_1$ on the square are
in~\citet{kuttler1972remarks} and~\citet{ferrero2005fourth}. Theorem~\ref{s6_highdim:thm:kappa}
adds a two-sided bound of $\kappa_d^2$ of the exact order $d^{-1/2}$ on the hypercube
(equivalently, $\kappa_d$ is of exact order $d^{-1/4}$), with an elementary proof (sine series,
the Cauchy--Schwarz inequality, partial fractions). In the literature searched for this work no
bound of $\kappa_d^2$ of the order $d^{-1/2}$ was found.
% src: research_highdim_bound/results/table_payne.tex; data/s6_highdim_torsion.json; refs.bib (bucur2011first, ferrero2005fourth, kuttler1972remarks, raulot2015sharp: annotes)
A posteriori bounds for network approximations exist in other norms and settings:
\citet{ernst2025certification} certify in the energy norm with constants not explicit in $d$;
\citet{fuhrer2025aposteriori} measure the boundary misfit in $H^{1/2}(\dOm)$ with mesh-dependent
constants; \citet{zeinhofer2025unified} work on $C^{1,1}$ domains with non-explicit constants;
in the generalisation bounds of~\citet{mishra2023generalization} the boundary misfit enters through
its square root with constants that contain norms of the network; and the bound
$\max\abs{v-\uex}\le\max_{\dOm}\abs{v-g}+\frac18\max\abs{\lap v+f}$ from the maximum principle,
explicit on the cube~\citep[Prop.~7]{mukherjee2026certification}, needs suprema, which the
Monte-Carlo losses do not estimate (on our networks its sampled value was 5.7 to 42 times the
bound of Corollary~\ref{s6_highdim:cor:bound}(c)). The closest form is in the error analysis of
residual minimisation by~\citet[App.~B, eq.~(B2)]{shin2023error}: $\norm{u}\le C(\norm{\lap
u}+\nbd{u})$, from the regularity theory of~\citet{bramble1970rayleigh}, linear in the $L^2$ norms
of the residual and of the boundary misfit as in Corollary~\ref{s6_highdim:cor:bound}(c), but on
smooth domains and with a constant that is not explicit. With Payne's constant a bound of this
form follows at once from the duality, so Corollary~\ref{s6_highdim:cor:bound}(c) is an application
of known theory with a better constant, not a new method. Imposing $v=g$
exactly~\citep{sukumar2022exact} would remove $B_{\mathrm{bd}}$ and the penalty bias; it was not
tested here. For the penalty bias, $O(1/\bet)$ is known on $C^{1,1}$ domains~\citep{muller2022},
and error bounds with extrapolation exist for the penalty finite-element
method~\citep{king1974penalty}; Proposition~\ref{sA_proofs:prop:exact} is a routine separation of
variables whose value lies in the numbers it gives on the cube.
% src: research_highdim_bound/results/decisions.json (H4: maxprinciple_over_L2bound: min 5.696, max 41.95); refs.bib (ernst2025certification, fuhrer2025aposteriori, zeinhofer2025unified, mishra2023generalization, mukherjee2026certification, shin2023error, sukumar2022exact, king1974penalty: annotes)

% -------------------------------------------------------------------------------------
\FloatBarrier
\subsection{Two cheap remedies at \texorpdfstring{$d=20$}{d = 20}}%
\label{s6_highdim:sec:remedies}

At $d=20$ and 4000 iterations neither method improves on the best affine approximation of
\Ptwo\ (Section~\ref{s6_highdim:sec:growth}); this is what is called the \emph{plateau} below: a
test error no smaller than the affine floor. Two cheap remedies suggest themselves, and neither
is new. The first, called \emph{presolve} here, writes the approximation as a fixed known part
plus a network, as in the trial solutions of \citet{lagaris1998}, with a cheap Galerkin solution
as the known part, as in Galerkin neural networks~\citep{ainsworth2021galerkin} and multi-level
networks~\citep{aldirany2024multilevel}: the fixed part is the exact minimiser of the method's own
population loss over affine functions, obtained from one $(d+1)\times(d+1)$ linear system
(Proposition~\ref{sA_proofs:prop:presolve}). The second, called \emph{lift3c}, replaces the input
$\x$ by the coordinate-wise Legendre features $P_1(t_i),P_2(t_i),P_3(t_i)$, $t_i=2x_i-1$, so that
the first layer has $3d$ inputs. Polynomial input layers trained with a PINN loss are known, for
high-frequency solutions: HOrderDNN~\citep{chang2022horderdnn}, whose first layer is a
tensor-product polynomial basis, and its coordinate-wise successor
K-HOrderDNN~\citep{zhang2025khorderdnn}, tested up to $d=10$ for PDEs and $d=50$ for function
fitting; related are the first layers of Chebyshev-basis Kolmogorov--Arnold
networks~\citep{shukla2024kan,mostajeran2025scaledcpikan} and feature-engineered or Fourier-feature
PINNs~\citep{fazliani2025safenet,wangwang2021}. Mapping the input to $[-1,1]^d$ is common practice
(the reference code of~\citet{raissi2019} does so), and the centred-input network $\mathcal N(2\x-1)$,
called \emph{lift1}, is reported as a baseline. That low-degree or low-frequency components are
learned first, which motivates both remedies, is known~\citep{rahaman2019,ghosh2022three,wangwang2021}.
What is measured here is whether either remedy removes the plateau on \Ptwo, and how each changes
the accuracy at $d=20$ on further problems, for both methods, at equal iterations and at equal CPU
time.
% src: refs.bib (chang2022horderdnn, zhang2025khorderdnn, raissi2019: annotes)

\paragraph{Design.} The test was pre-registered: hypotheses, problems, seeds 10 to 14, budgets,
metric and decision rule were fixed in a document whose hash and time were recorded in the
repository before the first counted run (no external registry was used). Both remedies had been
tried on \Pone\ and \Ptwo\ at $d=20$ in exploratory pilots with seeds 0 to 2 before the freeze, as
had the centred input (on \Pone\ at $d=20$, on \Ptwo\ up to $d=16$) and, for Deep Ritz on \Pone,
rep3, so the direction of their effects on these two problems was known. Only two problems were
held out, on which neither remedy had been tried before the test (a related ridge had been used in
an earlier pilot): the ridge \Pridge, $\uex=\cos(2s)$ with $s=d^{-1/2}\sum_i(2x_i-1)$, and
\Pcospair, $\uex=\sum_{k\le d/2}\cos(\pi x_{2k-1})\cos(\pi x_{2k})$. On \Ptwo\ the lift features
contain the solution up to $\sci{3.9}{-3}$, so \Ptwo\ is reported but excluded from the decision on
the lift. The metric is the relative $L^2$ error on the held-out test points, and the effect of an
arm is the seed-paired ratio $\rho=\errL(\text{plain})/\errL(\text{arm})$ ($\rho>1$: the arm is
more accurate). By the frozen rule a remedy \emph{helps} on a problem if, for both methods,
$\rho>1$ in at least 4 of 5 seeds with median $\rho\ge1.5$, at 4000 iterations and at equal CPU
time; it \emph{hurts} if for one method a median $\rho$ is at most $1/1.1$; it is a \emph{general
remedy} if it helps on \Pone\ and on \Pridge\ or \Pcospair\ and hurts on none of them. Two controls
were pre-registered, for Deep Ritz only: \emph{lift3u}, which spans the same functions as lift3c
but uses the uncentred feature $x_i$ in place of $P_1(t_i)$, and \emph{rep3}, the featureless input
$(t,t,t)$ with $t=2\x-1$; the centred input lift1 was added for both methods as an exploratory arm
after the frozen analysis. Three problems were added afterwards: under an addendum to the
pre-registration, the ridge \Paltridge, $\uex=\cos(\pi s'+1)$ with
$s'=d^{-1/2}\sum_i(-1)^{i+1}(2x_i-1)$ and $f=4\pi^2\uex$; and, in the descriptive second round of
the re-check, \Psinlin, $\uex=\sum_{k\le d/2}\sin(\pi x_{2k-1})\,x_{2k}$ with $f=\pi^2\uex$, and
\Pexpcos, $\uex=\sum_{k\le d/2}\ee^{x_{2k-1}}\cos x_{2k}$ with $f=0$. The full design, with the
weights, the equal-CPU budgets and the later additions, is in Supplementary
Section~\ref{S5:sec:remedies}.
% src: research_highdim_remedies/PREREG_C1.md (header: pilots with seeds 0-2; sections 2.2, 2.3, 7); notes/VERIFICATION.md (sec. 4.3, pilots before the freeze); research_highdim_remedies/PREREG_FREEZE.json; research_highdim_remedies/results/checks.json (C5_floors poisson/d20: exact_lift3); research_highdim_remedies/PREREG_C1_P5.md; research_highdim_remedies/results/check_p5_formulas.json; research_highdim_remedies/check/selftest2.json

\paragraph{Results.}
% src: research_highdim_remedies/results/decision.json (per_problem; general_lift3c); research_highdim_remedies/results/summary.json (P1-P4 presolve and lift3c); research_highdim_remedies/results/runs_main20.jsonl (ridge, ritz, seed 10); research_highdim_remedies/results/decision.json (presolve/P2 pinn: median 0.941, min 0.909, max 0.982, wins 0; centring_clause: lift3u median 0.523, 0 of 5; rep3_P1_descriptive: median 0.880, 0 of 5)
Presolve helps on \Pone\ only (median ratios 2.67 for Deep Ritz and 2.85 for the PINN, 5 of 5
seeds each; Table~\ref{s6_highdim:tab:remedies}). On \Ptwo\ it is neutral by the frozen rule
(median 1.00 for Deep Ritz) and for the PINN slightly worse in every seed (median 0.94, range 0.91
to 0.98); it changes nothing on \Pcospair, where the affine minimiser is zero
(Proposition~\ref{sA_proofs:prop:struct}), and hurts on the held-out ridge \Pridge\ (medians 0.66
and 0.79, no seed better than plain). On the four pre-registered problems lift3c is more accurate
than plain in 5 of 5 seeds for both methods, at 4000 iterations (median ratios 2.33 and 2.86 on
\Pone, 2.27 and 1.77 on \Pridge, 2.83 and 4.02 on \Pcospair, for Deep Ritz and the PINN) and at
equal CPU time, so by the frozen rule it is a general remedy on \Pone, \Pridge\ and \Pcospair.
The controls do not attribute this gain to the Legendre features on every problem
(Tables~\ref{sC_details:tab:controls} and~\ref{s6_highdim:tab:remedies5}(b)). On \Pone\ lift3u is
less accurate than plain in every seed (median 0.52), which triggered the pre-registered centring
clause, so no statement on the features is made there; no control reproduces the gain (rep3 0.88,
0 of 5 seeds better; lift1 0.92 for Deep Ritz, 2 of 5), and lift3c is 2.43 (Deep Ritz) and 1.61
(PINN) times more accurate than lift1, 5 of 5 seeds each. On \Pridge\ centring the input alone
matches the whole gain for Deep Ritz, and the PINN is more accurate with lift1 than with lift3c in
every seed (median 0.77). On \Pcospair\ lift3c also beats lift1 (2.23 and 2.39, 5 of 5), but for
Deep Ritz the featureless reparametrisation rep3 gains about as much as lift3c there, as on \Ptwo\
and \Pridge.
% src: research_highdim_remedies/results/attribution.json; research_highdim_remedies/results/summary_exploratory.json; research_highdim_remedies/results/summary.json (rep3, lift3c)
% src: research_highdim_remedies/PREREG_C1_P5.md; research_highdim_remedies/results/decision_p5.json; research_highdim_remedies/results/summary_p5.json; research_highdim_remedies/results/attribution.json (P5 pinn lift1_vs_plain 1.287, 5/5)
\Paltridge\ was run with the study's own code under the addendum, frozen before these runs; it is
not a blind test, since three seeds of the re-check had already shown the sign of the lift effect
on it. There the lift hurts (Deep Ritz median ratio 0.59 at 4000 iterations and 0.32 at equal CPU
time, 0 of 5 seeds better; PINN 0.88 and 0.72), and so does presolve (Deep Ritz 0.73, PINN 0.83),
whereas the centred input alone makes the PINN more accurate (1.29, 5 of 5). By the extended rule
of the addendum neither presolve nor lift3c is a general remedy on the five problems.
% src: data/s6_highdim_remedies_recheck.json (P6 lift1/lift3c 1.011; P7/ritz/d20: plain_err 0.02504-0.02532, presolve/4000 median 4.507); research_highdim_remedies/check/selftest2.json (P7_floor_affine 0.0224)
In the second round of the re-check (Deep Ritz, three seeds, a different implementation,
descriptive) the gain of the lift on \Psinlin\ is again matched by centring the input, and
\Pexpcos\ is on the plateau ($\errL=0.025$, affine floor 0.022), which presolve leaves (median
ratio 4.51).

\paragraph{What this shows and what it does not.} At $d=20$, with this network, optimiser,
schedule and these weights, the accuracy of both methods depends strongly on how the input enters
the first layer, in a way that differs between problems; no single mechanism (centring the input or
the polynomial features) accounts for it on every problem.
% src: research_highdim_remedies/results/summary.json (plain, d20); research_highdim_remedies/results/summary_p5.json (plain, d20); research_highdim_remedies/results/checks.json (C5_floors: test_affine); research_highdim_remedies/results/checks_p5.json (C5); data/s6_highdim_remedies_recheck.json (P6/ritz/d20 plain_err 0.0997-0.1155, P7/ritz/d20 plain_err 0.02504-0.02532); research_highdim_remedies/check/selftest2.json (P6_floor_affine 0.1716, P7_floor_affine 0.0224)
Of the five problems of the study only \Ptwo\ is on the plateau after 4000 iterations
(plain $\errL=0.126$ for Deep Ritz and $0.145$ for the PINN, against the affine floor $0.119$ on
the test set), and of the two of the re-check only \Pexpcos; on the others plain already ends
below the affine floor, so there the comparison is
one of accuracy at a fixed budget, not of escaping a plateau.
% src: research_highdim_remedies/results/summary.json (P2/d20: ritz lift3c, lift3u, rep3; pinn lift3c); research_highdim_remedies/results/summary_exploratory.json (P2/d20/ritz/lift1/4000: min 0.0699, max 0.0914; P2/d20/pinn/lift1/4000: min 0.0670, max 0.0796); research_highdim_remedies/results/checks.json (C5_floors poisson/d20: test_affine 0.1193, exact_lift3 3.9e-3); research_highdim_remedies/results/attribution.json and Table s6_highdim:tab:remedies5(b) (P2 lift1/lift3c: 4.22 Deep Ritz, 7.51 PINN, 5/5)
On \Ptwo\ presolve does not leave the plateau. lift3c does, for both methods and in every seed.
For Deep Ritz the gain cannot be attributed to its features, although they contain the solution of
\Ptwo\ up to $\sci{3.9}{-3}$: the featureless control rep3 is more accurate still (median ratio
8.10 against 6.59) and lift3u gains as much (6.71). The centred input alone (lift1, exploratory)
already ends below the affine floor in every seed for both methods ($\errL=0.070$ to $0.091$ for
Deep Ritz and $0.067$ to $0.080$ for the PINN, against 0.119), but with less gain than lift3c, which
is 4.22 times (Deep Ritz) and 7.51 times (PINN) more accurate than lift1 in median, 5 of 5 seeds
each (Table~\ref{s6_highdim:tab:remedies5}(b)); for the PINN, for which no featureless control was
run, the features thus add to centring on \Ptwo. At this budget the plateau of \Ptwo\ is therefore not robust to centring
the input, which is common practice. As a remedy the result is negative: neither remedy is general,
and neither is new as a method. What is added is a measurement of both, for the PINN and Deep Ritz
side by side, at equal iterations and at equal CPU time, on two problems held out from the pilots
and three added later, with two pre-registered controls for Deep Ritz (lift3u, rep3) and an
exploratory centred-input arm for both methods.

% -------------------------------------------------------------------------------------
\subsection{Comparison with published results}\label{s6_highdim:sec:limits}

% src: research_highdim/results/extensions.json (long, laplace, d = 10, 16 000 iterations: pinn 5.590e-3, 5.928e-3, 6.130e-3; ritz 6.830e-3, 7.438e-3, 7.568e-3); refs.bib (eyu2018, annote)
For \Pone\ at $d=10$, \citet{eyu2018} report a relative $L^2$ error of about 0.4 per cent after
50\,000 iterations with a 671-parameter network and $\bet=10^3$. Our mean errors at $d=10$ after
16\,000 iterations are 0.59 per cent (PINN) and 0.73 per cent (Deep Ritz), with best seeds at 0.56
and 0.68 per cent, for a different network, sampling and schedule. The orders of magnitude agree;
we do not claim to reproduce their number. For the Neumann example whose solution \Ptwo\ shares,
they report 1.3 per cent at $d=5$ and 1.9 per cent at $d=10$ with the Ritz method and no boundary
penalty; equation, boundary condition, network and budget differ from ours, so this is not a
like-for-like comparison with Table~\ref{s6_highdim:tab:p2}.
% src: refs.bib (chen2020comparison, annote: abstract of arXiv:2005.04554, checked 2026-10-01)
\citet{chen2020comparison} set the deep Galerkin method~\citep{sirignano2018}, whose loss is the
residual least-squares loss of a PINN, against the Deep Ritz method on elliptic problems with
boundary conditions imposed by penalty. They report that Deep Ritz can outperform the residual
method even for smooth solutions, with a clear dependence on the dimension, that the residual
method can outperform Deep Ritz for solutions of low regularity, and that imposing the boundary
condition exactly, where that is possible, gives a better solution and an easier training. Our
finding that the PINN is the more accurate method per iteration for $d\le10$ on two smooth
problems is therefore not a general ranking, and it does not contradict theirs: networks,
sampling, penalty weights and budgets differ, and our own ranking changes with what is counted
and with the budget. What the two studies have in common is that neither loss dominates the
other.
